\documentclass[tikz]{standalone}
\usepackage{ctex}
\usepackage{tikz}
\usepackage{pgfplots}
\pgfplotsset{compat=1.18}   % 使用较新版本语法
\usetikzlibrary{positioning, arrows.meta, intersections, calc, fit, backgrounds}
\usepgfplotslibrary{fillbetween}  % 用于填充区域

\begin{document}
\begin{tikzpicture}[->,>=stealth,shorten >=1pt,auto,
        thick,
        branch/.style={rectangle,draw,rounded corners=2pt,fill=blue!5,font=\sffamily\small,minimum width=1.8cm,minimum height=0.6cm},
        main node/.style={circle,draw,font=\sffamily\small}
    ]

    % 前一层输出
    \node[main node] (prev) {前一层};

    % 四个并行分支
    \node[branch] (b1) at (3, 1.8) {$1\times1$ 卷积};
    \node[branch] (b2) at (3, 0.6) {$3\times3$ 卷积};
    \node[branch] (b3) at (3, -0.6) {$5\times5$ 卷积};
    \node[branch] (b4) at (3, -1.8) {$3\times3$ 最大池化};

    % 拼接节点
    \node[main node] (concat) at (6, 0) {$\|$};

    % 下一层
    \node[main node] (next) at (8.5, 0) {下一层};

    % 连线：前一层到各分支
    \path[every node/.style={font=\sffamily\small}]
    (prev) edge[out=0, in=180] (b1)
    (prev) edge[out=0, in=180] (b2)
    (prev) edge[out=0, in=180] (b3)
    (prev) edge[out=0, in=180] (b4);

    % 各分支到拼接
    \path[every node/.style={font=\sffamily\small}]
    (b1) edge (concat)
    (b2) edge (concat)
    (b3) edge (concat)
    (b4) edge (concat);

    % 拼接到下一层
    \draw[->] (concat) -- (next);

    % 说明性文字
    \node[above right=0.2cm and 0.1cm of concat, font=\sffamily\scriptsize] {拼接};
\end{tikzpicture}
\end{document}